\documentclass[11pt]{article}
\usepackage[a4paper,margin=1.9cm]{geometry}
\usepackage[T1]{fontenc}
\usepackage{textcomp}
\usepackage[spanish]{babel}
\usepackage{amsmath,amssymb}
\usepackage{enumitem}
\usepackage{tikz}
\usetikzlibrary{calc,arrows.meta,patterns,decorations.pathmorphing}
\usepackage{xcolor}
\usepackage{multicol}
\usepackage{parskip}
\usepackage{lmodern}
\renewcommand{\familydefault}{\sfdefault}

\definecolor{unalblue}{RGB}{31,73,124}

\newlist{opciones}{enumerate}{1}
\setlist[opciones]{label=(\alph*),itemsep=1pt,topsep=2pt,leftmargin=1.8em,font=\normalfont}
\newcommand{\ptsbox}[1]{\fbox{\footnotesize #1}~}

\newcommand{\vect}[2]{\begin{pmatrix}#1\\#2\end{pmatrix}}
\newcommand{\vectiii}[3]{\begin{pmatrix}#1\\#2\\#3\end{pmatrix}}

\newcommand{\examheader}{%
  \noindent
  \begin{minipage}[t]{0.6\linewidth}
    {\small\bfseries Geometría Vectorial y Analítica --- Parcial 2}\\
    {\small Semestre 2014--01}
  \end{minipage}%
  \hfill
  \begin{minipage}[t]{0.35\linewidth}
    \raggedleft\small Escuela de Matemáticas. Universidad Nacional de Colombia, Sede Medellín.
  \end{minipage}\\[2pt]
  \rule{\linewidth}{0.6pt}
}
\newcommand{\examfooter}{%
  \vfill
  \rule{\linewidth}{0.4pt}\\[1pt]
  {\scriptsize\textit{Transcripción digital --- MathUNAL}\hfill\textcolor{unalblue}{\textbf{\thepage}}}
}
\pagestyle{empty}

\begin{document}

\examheader

\begin{center}
{\bfseries Segundo examen parcial de Geometría Vectorial}\\
5 de mayo de 2014

\vspace{4pt}
\textbf{Puntaje. Solo para uso oficial}

\vspace{2pt}
\begin{tabular}{|c|c|c|c|c|c|}
\hline
1 & 2 & 3 & 4 & 5 & TOTAL\\
\hline
 & & & & & \\
\hline
\end{tabular}
\end{center}

\textbf{Instrucciones:} No está permitido el uso de calculadora ni de otros aparatos electrónicos. Solucione las preguntas en los espacios en blanco asignados a cada una de ellas. La última página (página 6) sólo puede usarse para operaciones en borrador; nada de lo escrito en esta página se tendrá en cuenta como parte de la solución del examen. Duración del examen: \textbf{1 hora y 50 minutos}. No está permitido sacar hojas en blanco ni ningún tipo de apuntes durante el examen.

\vspace{10pt}

\begin{enumerate}[leftmargin=1.6em]

\item Sean $\mathcal{L}$ la recta con ecuación $x+2y=0$, $P$ la transformación proyección ortogonal sobre la recta $\mathcal{L}$ y $S$ la transformación reflexión con respecto a la recta $\mathcal{L}$.

\begin{opciones}
\item \ptsbox{4 pts} Complete cada expresión para obtener una afirmación verdadera:

\textbf{(i)} \ptsbox{4 pts} La matriz de la transformación $P$ es
\[
m(P) = \begin{pmatrix} \phantom{-\dfrac{0}{0}} & \phantom{-\dfrac{0}{0}} \\[6pt] \phantom{-\dfrac{0}{0}} & \phantom{-\dfrac{0}{0}} \end{pmatrix}
\]
\textbf{(ii)} \ptsbox{4 pts} La matriz de la transformación $S$ es
\[
m(S) = \begin{pmatrix} \phantom{-\dfrac{0}{0}} & \phantom{-\dfrac{0}{0}} \\[6pt] \phantom{-\dfrac{0}{0}} & \phantom{-\dfrac{0}{0}} \end{pmatrix}
\]
\textbf{(iii)} \ptsbox{4 pts} Si $\mathcal{L}'$ es la recta con ecuación $y=2x$, la imagen de $\mathcal{L}'$, bajo la transformación $P$ es:

\vspace{20pt}
\textbf{(iv)} \ptsbox{4 pts} La imagen de la recta $\mathcal{L}'$, definida en (iii), bajo la transformación $S$ es la recta con ecuación:

\vspace{20pt}

\item \ptsbox{12 pts} Escriba, dentro del paréntesis, V si la afirmación es verdadera o F si es falsa:

\textbf{(i)} \ptsbox{3 pts} La transformación $P$ es invertible. \hspace{1cm} ( \phantom{x} )\\[4pt]
\textbf{(ii)} \ptsbox{3 pts} La imagen de $\mathbb{R}^2$ bajo $S\circ P$ es $\mathcal{L}$. \hspace{1cm} ( \phantom{x} )\\[4pt]
\textbf{(iii)} \ptsbox{3 pts} La imagen de $\mathcal{L}$ bajo $R_{\pi/2}\circ P$ es la recta con ecuación $2x-y=0$. \hspace{1cm} ( \phantom{x} )\\[4pt]
\textbf{(iv)} \ptsbox{3 pts} $P$ es un movimiento euclidiano. \hspace{1cm} ( \phantom{x} )
\end{opciones}

\item Sea $T$ una transformación lineal para la cual
\[
T\vect{1}{2} = \vect{2}{4} \qquad \text{y} \qquad T\vect{2}{1} = \vect{0}{0},
\]
y sea $A$ la matriz de $T$.

\begin{opciones}
\item \ptsbox{6 pts} ¿Es $T$ invertible? Justifique su respuesta.

\vspace{40pt}

\item \ptsbox{6 pts} ¿Cuáles son los valores propios de $A$?

\vspace{40pt}

\item \ptsbox{8 pts} Encuentre la matriz $A$.

\vspace{60pt}
\end{opciones}

\item
\begin{opciones}
\item \ptsbox{12 pts} Sea $X_1 = \vect{1}{-3}$. Halle un vector $X_2$ tal que el ángulo de $X_1$ a $X_2$ es $30$° y el área del triángulo $\mathcal{T}$ con vértices $\vect{0}{0}$, $X_1$ y $X_2$ es 5 unidades cuadradas.

\vspace{80pt}

\item \ptsbox{12 pts} Si $T$ es la transformación lineal definida por $T\vect{x}{y} = \vect{2x+y}{x-2y}$, halle el área de la imagen bajo $T$ del triángulo $\mathcal{T}$ del literal (a).

\vspace{60pt}
\end{opciones}

\newpage
\examheader

\item
\begin{opciones}
\item \ptsbox{3 pts} Escriba la definición de \emph{Movimiento euclidiano}.

\vspace{30pt}

\item \ptsbox{3 pts} Sea $U$ un vector no nulo del plano y sea $T_U : \mathbb{R}^2 \to \mathbb{R}^2$ la transformación dada por $T_U(X) = X+U$. Demuestre que $T_U$ es un movimiento euclidiano.

\vspace{50pt}

\item \ptsbox{3 pts} Escriba la definición de \emph{Matriz ortogonal}.

\vspace{30pt}

\item \ptsbox{3 pts} Demuestre que si $Q$ es una matriz ortogonal entonces $\det(Q)=1$ o $\det(Q)=-1$.

\vspace{50pt}
\end{opciones}

\item \ptsbox{16 pts} Utilice un movimiento euclidiano (una rotación) para transformar la ecuación
\[
16x^2 - 24xy + 9y^2 - 30x - 40y = 0
\]
en una ecuación de la forma $A'(x')^2 + B'(y')^2 + D'x' + E'y' + F' = 0$.

\vspace{100pt}

\end{enumerate}

\examfooter
\end{document}
